\section{Experimental setup and event selection}\label{Sec:ALICE}

A detailed description of the ALICE apparatus in its Run 1 and 2 configuration, as well as its performance, can be found in Refs.~\cite{ALICE1,ALICE2}. This section will briefly describe the main ALICE detectors used for the event selection, \SOPT determination, and extraction of particle spectra. All radii in the following are given as distances from the beam axis.

The detectors of the ALICE apparatus can be grouped as follows: the central barrel at midrapidity, the muon arm at forward rapidity, and the forward global detectors. To be able to efficiently trigger on inelastic \pp collisions, ALICE employs two forward scintillator arrays, V0A and V0C, with a pseudorapidity coverage of $2.8 < \eta < 5.1$ and $-3.7 < \eta < -1.7$,
respectively. 


Both \SOPT and particle yields are determined by using tracks in the central barrel. The main detectors for tracking in the central barrel are the Inner Tracking System (ITS) and the Time Projection Chamber (TPC). The ITS detector is composed of six layers of cylindrical silicon detectors with full azimuthal acceptance, where the radius of the innermost (outermost) layer is $3.9$\,cm ($43$\,cm). The innermost layers consist of two arrays of hybrid silicon pixel detectors (SPD), whose fine granularity provides high precision tracking closest to the primary vertex (PV). The SPD is also used to reconstruct tracklets, short two-point track segments covering the pseudorapidity region $\etarange{1.4}$. The tracklets provide both efficient primary vertex determination and a precise estimate of the charged-particle multiplicity. The remaining layers of the ITS are only used for tracking in the analyses presented here. The TPC is the primary tracking device in ALICE that provides a full three-dimensional trajectory for  each track. It is a large cylindrical detector that surrounds the ITS detector with an inner and outer radii of 85\,cm and 250\,cm, respectively.  It has full azimuthal acceptance and a pseudorapidity coverage of $-0.9< \eta <0.9$ for full-length tracks. The ITS and TPC are situated inside the solenoidal L3 magnet with a uniform magnetic field of 0.5\,T. For global tracks, where the full information of the ITS and TPC are used together, a momentum resolution of 1--10\,\% is achieved for momenta ranging from 0.05 to \gevc{100}.

The ALICE apparatus utilizes a broad range of different particle identification techniques to identify the mass of each analyzed particle. For weakly-decaying $V^0$ (\kzero, \lmb) and Cascades ($\Xi$), the TPC tracking is able to provide topological track matching of the decay products, where particles are then identified through peaks in invariant mass distributions. For the other particles, the PID is provided by the specific energy loss, d$E$/d$x$, measured in the TPC and the particle's velocity (given a measured momentum in the TPC) is provided by the Time-of-flight (TOF) detector. The TOF consists of Multi-Gap Resistive Plate Chambers (MRPC), which are used for the particle identification by measuring the total time of flight of the identified hadrons. It is located at about 3.7\,m from the interaction point, with a full azimuthal acceptance and has a pseudorapidity coverage of $-0.9< \eta <0.9$. Details of the particle identification procedure are given in Sec.~\ref{Sec:Details}.
